\section*{Capítulo \ref{ch:b2:palaeoenvironments} --- Fósseis e paleoambientes}

\begin{solution}{exo:b2:palaeoenvironments:1}
(a) $112/135 = 0.830$ e $112/150 = 0.747$. (b) As conchas mortas no fundo se acumulam durante anos
ou mais: incluem espécies que viveram ali antes, ou raramente, ou que foram trazidas, de modo que a
\omterm{def:b2:palaeoenvironments:assemblages}{assembleia de morte} faz a média de muitas gerações da comunidade. (c) Os grupos de corpo mole, em
particular os detritívoros e os carnívoros, como os vermes: sem esqueleto, decompõem-se antes do soterramento.
\end{solution}

\begin{solution}{exo:b2:palaeoenvironments:2}
(a) $5 \times 165 = \qty{825}{kyr}$, cerca de \qty{0.8}{Myr}, dentro dos 0.7 a 1 milhão de anos.
(b) $183.73 - 182.77 = \qty{0.96}{Myr}$, e $0.96/0.165 = 5.8$: cerca de seis horizontes. (c) Uma
linhagem muda em muito menos de um milhão de anos; só os horizontes separam as formas
sucessivas, enquanto uma zona as agrupa.
\end{solution}

\begin{solution}{exo:b2:palaeoenvironments:3}
(a) \omterm{def:b2:palaeoenvironments:ichnofacies}{Icnofácies} Nereites: o fundo de mar profundo de um sistema turbidítico, além da quebra da plataforma, onde
o alimento é escasso. (b) \omterm{def:b2:palaeoenvironments:ichnofacies}{Icnofácies} Skolithos: estirâncio ou antepraia. (c) \omterm{def:b2:palaeoenvironments:ichnofacies}{Icnofácies} Cruziana:
costa afora, entre os níveis de base das ondas de tempo bom e de tempestade. (d) \omterm{def:b2:palaeoenvironments:ichnofacies}{Icnofácies} Zoophycos: no
Mesozoico, a plataforma inferior ou o talude, abaixo do nível de base das ondas de tempestade.
\end{solution}

\begin{solution}{exo:b2:palaeoenvironments:4}
(a) A camada A contém grupos estenoalinos: água do mar normal. A camada B, pobre em espécies mas rica em
indivíduos, com restos vegetais e nenhum dos grupos estenoalinos, sugere água salobra, por
exemplo um estuário. (b) Os grupos ausentes são vários, comuns e bem preservados em camadas marinhas
da mesma idade: sua ausência conjunta é um padrão, enquanto um fóssil pode faltar por acaso ou
por má preservação. (c) O sedimento (canais de maré ou fluviais, drapeamentos de lama, raízes), os
\omterm{def:b2:palaeoenvironments:fossil-kinds}{icnofósseis} (uma \omterm{def:b2:palaeoenvironments:ichnofacies}{icnofácies} Cruziana com poucas espécies), e a passagem lateral da camada B a
depósitos marinhos ou fluviais.
\end{solution}

\begin{solution}{exo:b2:palaeoenvironments:5}
(a) $P = 1/2$ quando $kh/s = \ln 2$: $s = 0.05 \times 0.1/0.693 = \qty{0.0072}{m.yr^{-1}}$,
\qty{7.2}{mm} por ano. (b) Reduzir $s$ à metade dobra o expoente: $P = (1/2)^2 = 0.25$, a metade.
(c) $k$ cai à metade: $P = e^{-2.5} = 0.082$ em vez de $e^{-5} = 0.0067$, $e^{2.5} = 12$ vezes mais
provável. (d) As conchas espessas e robustas ficam super-representadas, e tanto mais quanto mais lento o soterramento.
\end{solution}

\begin{solution}{exo:b2:palaeoenvironments:6}
(a) O limite superior é $183.73 + 0.35 = \qty{184.08}{Ma}$: os \qty{184.2}{Ma} da carta ficam
fora dele, por \qty{0.12}{Myr}. (b) $183.73 - 182.77 = \qty{0.96}{Myr}$. (c) $6/0.96 =
\qty{6.25}{m}$ por milhão de anos. (d) Uma biozona é definida por fósseis, cujas aparições e
desaparições são quase simultâneas em áreas amplas; a espessura de rocha depositada durante
esse tempo depende da taxa local de acumulação.
\end{solution}

\begin{solution}{exo:b2:palaeoenvironments:7}
(a) $0.49\,e^{-1046/3704} = 0.369$ e $0.49\,e^{-1166/3704} = 0.358$; média $0.364$. (b) $120
\times (1 - 0.364) = \qty{76.3}{m}$. (c) $76.3/(1 - 0.49) = \qty{149.6}{m}$. (d) O fator
supõe que toda a camada restaurada tenha a porosidade $\varphi_0$ da superfície; mas uma camada
de \qty{150}{m} de espessura na superfície é menos porosa em sua base que em seu topo, de modo que os mesmos grãos
ocupam menos de \qty{149.6}{m}.
\end{solution}

\begin{solution}{exo:b2:palaeoenvironments:8}
(a) Um volume unitário contém $\varphi$ de água e $1 - \varphi$ de grãos: $\rho = \varphi
\rho_{\text{water}} + (1 - \varphi)\rho_{\text{grain}}$, o que dá a fórmula; $(2710 -
2573)/(2710 - 1000) = 0.080$. (b) $0.51\,e^{-956/4545} = 0.413$. (c) A cimentação: a calcita
precipitada nos poros, que os carbonatos sofrem além da compactação, elimina muito mais porosidade
que o soterramento sozinho.
\end{solution}

\begin{solution}{exo:b2:palaeoenvironments:9}
(a) Turbiditos de mar profundo (\omterm{def:b2:palaeoenvironments:ichnofacies}{icnofácies} Nereites); um fundo jurássico da plataforma inferior ou do talude
(Zoophycos); um fundo costa afora entre os níveis de base das ondas (Cruziana); uma antepraia (Skolithos); um
pântano em terra. (b) A água ficou mais rasa para cima: a costa deslocou-se em direção ao mar. (c) Uma descida do nível do mar,
ou um aporte de sedimento mais rápido que a subsidência, como diante de um delta em crescimento; uma descida
seria encontrada na mesma época em outras bacias, um delta seria local e seu sedimento viria
de um lado só. (d) Numa sucessão concordante os ambientes vizinhos seguem-se uns aos outros; um
salto por cima dos intermediários aponta para uma lacuna no tempo ou para uma falha.
\end{solution}

\begin{solution}{exo:b2:palaeoenvironments:10}
(a) O tempo até o soterramento é $h/s$ = 100, 10 e \qty{1}{yr}. Corpo mole: $e^{-100} \approx
\num{4e-44}$, $e^{-10} = \num{4.5e-5}$, $e^{-1} = 0.37$. Concha ($k = \qty{0.01}{yr^{-1}}$):
$e^{-1} = 0.37$, $e^{-0.1} = 0.90$, $e^{-0.01} = 0.99$. (b) $s = kh/\ln 100 = 0.1/4.605 =
\qty{0.022}{m.yr^{-1}}$, \qty{22}{mm} por ano. (c) Dificilmente, durante meses; mas uma camada de tempestade, um
turbidito ou uma queda de cinzas soterra o fundo em horas. (d) A taxa constante de sedimentação: o
tempo até o soterramento é então muito mais curto que $h/s$ com a taxa média, e $P = e^{-kT}$, sendo $T$
o tempo real até o soterramento, tende a um quando $T$ é muito menor que $1/k$.
\end{solution}

\begin{solution}{exo:b2:palaeoenvironments:11}
(a) Uma camada fina $\mathrm dz$ à profundidade $z$ contém os grãos $(1 - \varphi(z))\,\mathrm dz$;
na superfície eles ocupam $\mathrm dz_0$ com $(1 - \varphi_0)\,\mathrm dz_0 = (1 -
\varphi(z))\,\mathrm dz$. (b) $\varphi = 0.63\,e^{-816/1960} = 0.415$; $(1 - 0.415)/0.37 =
1.581$, portanto \qty{158.1}{m}. (c) \qty{9.5}{m} a mais: como no \cref{exo:b2:palaeoenvironments:7},
a camada restaurada não é uniformemente tão porosa quanto sua superfície. (d) A taxa de grãos,
$h(1-\varphi)/\Delta t$, usa apenas a espessura e a porosidade atuais, ambas medidas.
\end{solution}

\begin{solution}{exo:b2:palaeoenvironments:12}
(a) $9.5/0.165 = 57.6$ horizontes. (b) $180/57.6 = \qty{3.1}{m}$ de calcário hoje, $338/57.6 =
\qty{5.9}{m}$ de sedimento fresco. (c) No máximo \qty{1.5}{m}: $180/1.5 = 120$ intervalos, 121
amostras. (d) Interrupções na deposição, camadas onde as amonites estão ausentes, e horizontes condensados em
alguns centímetros onde o sedimento chegava devagar.
\end{solution}

\begin{solution}{pb:b2:palaeoenvironments:1}
\textbf{1.} $866 \sin \ang{12} = 866 \times 0.2079 = \qty{180.0}{m}$.
\textbf{2.} $481 \times 0.2079 = \qty{100.0}{m}$.
\textbf{3.} O arenito, por superposição: os flancos mergulham apenas \ang{12}, de modo que as camadas não estão
invertidas e a mais baixa é a mais antiga.
\textbf{4.} $866 - 766 = \qty{100}{m}$ e $1046 - 866 = \qty{180}{m}$.
\textbf{5.} Pelo menos \qty{766}{m} ali; seu topo original foi erodido, de modo que sua espessura completa é
desconhecida.
\textbf{6.} 0.70 na superfície; $0.70\,e^{-766/1408} = 0.406$ a \qty{766}{m}.
\textbf{7.} São grupos estenoalinos: a água tinha a salinidade do mar aberto.
\textbf{8.} Uma \omterm{def:b2:palaeoenvironments:ichnofacies}{icnofácies} Cruziana: a zona costa afora, entre os níveis de base das ondas de tempo bom e de
tempestade.
\textbf{9.} Tempestades: suas ondas atingiam o fundo, quebravam as conchas e as depositavam em camadas
com bases erosivas.
\textbf{10.} A água de fundo era pobre em oxigênio: poucos escavadores, lâminas conservadas, pirita formada.
\textbf{11.} A \omterm{def:b2:palaeoenvironments:ichnofacies}{icnofácies} Zoophycos, no Mesozoico a plataforma inferior ou o talude, abaixo do
nível de base das ondas de tempestade; oxigênio suficiente, por vezes, para os escavadores.
\textbf{12.} Não: as amonites nadavam na água acima; suas conchas afundavam ou derivavam até o fundo
depois da morte e formam uma \omterm{def:b2:palaeoenvironments:assemblages}{assembleia de morte}.
\textbf{13.} $184.2 - 174.7 = \qty{9.5}{Myr}$.
\textbf{14.} $180/9.5 = \qty{18.9}{m}$ por milhão de anos.
\textbf{15.} $(2710 - 2573)/(2710 - 1000) = 0.080$.
\textbf{16.} $180 \times (1 - 0.080)/(1 - 0.51) = \qty{338}{m}$.
\textbf{17.} $338/9.5 = \qty{35.6}{m}$ por milhão de anos.
\textbf{18.} $180 \times 0.92/9.5 = \qty{17.4}{m}$ por milhão de anos.
\textbf{19.} O folhelho deu 11.3 e \qty{4.4}{m} por milhão de anos: a taxa de grãos caiu
cerca de quatro vezes. O carbonato do baixio e da plataforma deixou de chegar à medida que o fundo
se aprofundava, e só uma chuva lenta de lama o atingia.
\textbf{20.} Avançou: rios, depois um baixio, um fundo costa afora, um fundo abaixo do nível de base das ondas de
tempestade, depois o mar aberto do giz. A coluna registra uma transgressão.
\textbf{21.} $57 - 41 = \qty{16}{m}$ durante o calcário, $86 - 57 = \qty{29}{m}$ durante o folhelho,
\qty{45}{m} no total.
\textbf{22.} $(60 - 5) + 338 + 148.6 = \qty{542}{m}$, em $184.2 - 161.5 = \qty{22.7}{Myr}$:
\qty{23.9}{m} por milhão de anos.
\textbf{23.} $45/542 = 0.083$: cerca de \qty{8}{\percent}, e menos se a carta de ciclos exagera
a amplitude. O resto veio da subsidência do fundo.
\textbf{24.} A lama calcária do calcário de Varden acumulou-se a cerca de
\textbf{\qty[text-series-to-math=true]{36}{m}} por milhão de anos quando foi depositada.
\end{solution}
